\documentclass{article}

% Pass 'final' to the neurips_2023 package to remove line numbers if needed
\usepackage[final]{neurips_2023}

\usepackage{graphicx}

\usepackage{enumitem}

% \usepackage{placeins}

\usepackage[utf8]{inputenc} % allow utf-8 input
\usepackage[T1]{fontenc}    % use 8-bit T1 fonts
\usepackage{hyperref}       % hyperlinks
\usepackage{url}            % simple URL typesetting
\usepackage{booktabs}       % professional-quality tables
\usepackage{amsfonts}       % blackboard math symbols
\usepackage{nicefrac}       % compact symbols for 1/2, etc.
\usepackage{microtype}      % microtypography
\usepackage{xcolor}         % colors

\usepackage{hyperref}
\usepackage{cleveref}

\title{Summarization: Path Planning for Manipulation using Experience-driven Random Trees}

% Paper BibTex
% @article{pairet2021path,
%   title={Path planning for manipulation using experience-driven random trees},
%   author={Pairet, {\`E}ric and Chamzas, Constantinos and Petillot, Yvan and Kavraki, Lydia E},
%   journal={IEEE Robotics and Automation Letters},
%   volume={6},
%   number={2},
%   pages={3295--3302},
%   year={2021},
%   publisher={IEEE}
% }

\author{
  Summarized by Xiaonan (Nice) Wang \thanks{wangxiaonannice@gmail.com} \\
  % Paper fr last-names (year)
  Summarization\thanks{\textbf{Disclaimer:} This summarization is for research and learning purposes. It represents a personal interpretation and may contain inaccuracies. Feedback or corrections via email are highly appreciated.} Generated based on Paper from Pairet et al. (2021) \\
  Topic: \\
}

\begin{document}

\newcommand{\repeatfootnote}[1]{\textsuperscript{\ref{#1}}}

\maketitle

\begin{abstract}
This document summarizes the core contributions and methodology of the paper "Path Planning for Manipulation using Experience-driven Random Trees, \cite{pairet2021path}", focusing on its' main ideas and the core blocks.
\end{abstract}

\section{Overview: Core Questions and Answers}
                           
\paragraph{(1) What is the problem?} 
\begin{itemize}
  \item Motion Planning from \textbf{Prior Experiences}
  \item \textbf{Generalizability:} Generalize prior knowledge (experiences) even to \textbf{notably different} task instances.
\end{itemize}

\paragraph{(2) Why need to solve this problem?}
\begin{itemize}
  \item Robots planning from \textit{scratch} leads to \textit{unnecessarily long planning times}.
  \item While leveraging prior experiences can \textbf{speed up} the motion planning process, existing learning-based and experience-based methods \textit{fail} or \textit{generalize poorly} when the new task \textit{differs significantly} from the \textit{stored} experiences (i.e. \textit{non-experienced regions}).
  \item Making \textbf{prior} knowledge \textbf{truly reusable} across \textbf{diverse task variations} remains an open challenge.
\end{itemize}

\paragraph{(3) How is it different from prev.?}
\begin{itemize}
  \item Previous experience-based \textit{path-centric} planners (like \textit{Lightning} or \textit{Thunder}) use prior experiences \textit{``rigidly''} by recalling \textit{exact existed motions} and simply repairing disconnected states, which requires \textit{a vast library of closely related experiences}.
  \item In contrast, this approach treats experiences as \textbf{``decomposable''} and \textbf{``malleable''}. Instead of \textit{copying} exact paths, it breaks a prior path into smaller parts (\textbf{``micro-experiences''}) and \textbf{semi-randomly morphs} them using affine transformations (\textbf{shear} and \textbf{shift}) to explore new situations.
\end{itemize}

\paragraph{(4) Why is it better than prev.? (Advantages)} 
\begin{itemize}
  \item The proposed method can successfully solve \textbf{highly novel} task instances using just a \textbf{single} prior path experience, whereas \textit{state-of-the-art} planners fail even with \textit{vast repertoires of experiences}.
  \item It \textbf{significantly outperforms} existing methods in both \textbf{success rate} and \textbf{planning time}, particularly in challenging scenarios involving \textbf{new obstacles} or \textbf{narrower geometries}.
\end{itemize}

\paragraph{(5) What is the approach itself?}
\begin{itemize}
  \item Port-Hamiltonian neural ODE network on a Lie group
  \begin{itemize}
    \item 4 separate neural networks to \textbf{approximate physical components}:
    \begin{itemize}
      \item kinetic energy $T$
      \item potential energy $V$
      \item energy dissipation matrix $D$
      \item input matrix $B$
    \end{itemize}
    \item The networks dictate the continuous-time dynamics within an ODE solver to minimize the errors between predicted and true states.
  \end{itemize}
  \item Unified Control Policy: It applies an energy-shaping and damping-injection strategy to reshape the learned system's energy landscape, so its global minimum aligns with the desired tracking trajectory.
\end{itemize}

\paragraph{(6) What are the applications of it?} 
\begin{itemize}
  \item The approach (designed \textit{dynamics models}) is applicable to rigid-body \textbf{mobile} robots, such as \textbf{unmanned ground vehicles (UGVs)}, \textbf{unmanned aerial vehicles (UAVs)} and \textbf{unmanned underwater vehicles (UUVs)}.
  \item Experiments both in \textit{simulation} and \textit{\textbf{real} PX4 quadrotors}.
\end{itemize}

\section{The Structure}

\paragraph{Summarized Block-Diagram} 

% The summarized block-diagram of \textit{Port-Hamiltonian ODE} see \cref{fig:port_hamiltonian_ode}\footnote{For efficiency, this diagram was initially hand-drawn on paper and then converted into its current digital version using Nano Banana Pro.}.

% \begin{figure}[tb]
%   \centering
%   \includegraphics[width=0.95\textwidth]{fig/port_hamiltonian_ode.png}
%   \caption{Summarized Block-Diagram of Port-Hamiltonian ODE}
%   \label{fig:port_hamiltonian_ode}
% \end{figure}

\paragraph{Original in Paper}

% The original block-diagram see \cref{fig:original}.

% \begin{figure}[tb]
%   \centering
%   \includegraphics[width=0.95\textwidth]{fig/orig.png}
%   \caption{Original Block-Diagram in Paper}
%   \label{fig:original}
% \end{figure}

\section{Open Questions}

% \begin{enumerate}
%     \item \textbf{Generalizing to Highly Stochastic Environments:}
%     In highly stochastic environments, such as rocky terrains or slippery surfaces, complex maneuvers are often essential for survival. Does forcibly compressing the fractal dimension risk compromising the agent's performance in high-difficulty settings? Furthermore, is there an ``optimal dimension'' that strikes a balance between guaranteed verifiability and environmental adaptability?
%     \item \textbf{Beyond the Model-Free RL:} 
%     Beyond the field of Model-Free RL, does the concept of Fractal Regularization still hold significant value in legged locomotion based on recent popular approaches?
% \end{enumerate}

\bibliographystyle{abbrvnat}
\bibliography{main}
\end{document}
